% ---------------------------------------------------------------
%  DON YEU CAU CONG NHAN SANG KIEN  (Mau 3)
% ---------------------------------------------------------------
\thispagestyle{empty}
\begin{center}
{\bfseries CỘNG HÒA XÃ HỘI CHỦ NGHĨA VIỆT NAM}\\
{\bfseries Độc lập -- Tự do -- Hạnh phúc}\\[4pt]
\rule{5cm}{0.5pt}\\[1cm]
{\Large\bfseries ĐƠN YÊU CẦU CÔNG NHẬN SÁNG KIẾN}\\[1cm]
\end{center}

\noindent\textbf{Kính gửi:} Hội đồng Sáng kiến \dotfill
\ghichu{Điền đúng tên hội đồng ghi trong công văn hướng dẫn của Sở
(thường là ``Hội đồng Sáng kiến ngành Giáo dục và Đào tạo tỉnh \dots'').}

\medskip
\noindent Tôi ghi tên dưới đây:

\begin{center}
\begin{tabular}{|c|p{2.6cm}|p{2cm}|p{2.4cm}|p{1.8cm}|p{2cm}|p{1.5cm}|}
\hline
\textbf{TT} & \textbf{Họ và tên} & \textbf{Ngày sinh} & \textbf{Nơi công tác} &
\textbf{Chức danh} & \textbf{Trình độ chuyên môn} & \textbf{Tỷ lệ đóng góp} \\ \hline
1 & \TACGIA & \NAMSINH & \DONVI & \CHUCVU & \TRINHDO & 100\% \\ \hline
\end{tabular}
\end{center}

\ghichu{Điền ngày sinh, trình độ chuyên môn (Cử nhân/Thạc sĩ Toán học) vào file
\texttt{skkn.tex} ở phần ``THÔNG TIN CẦN ĐIỀN'', cả bài sẽ tự cập nhật.}

\bigskip
\noindent\textbf{1. Là tác giả đề nghị xét công nhận sáng kiến:} \TENSANGKIEN

\noindent\textbf{2. Chủ đầu tư tạo ra sáng kiến:} \TACGIA\ (tác giả tự đầu tư
kinh phí thuê máy chủ, tên miền và tự thực hiện toàn bộ phần lập trình).

\noindent\textbf{3. Lĩnh vực áp dụng sáng kiến:} \LINHVUC. Sáng kiến được áp
dụng trong khâu biên soạn đề kiểm tra và tổ chức kiểm tra, đánh giá môn Toán
cấp trung học phổ thông, phục vụ đồng thời giáo viên (lấy đề, giao bài, theo
dõi kết quả lớp) và học sinh (tự luyện tập, làm bài trực tuyến).

\noindent\textbf{4. Ngày sáng kiến được áp dụng lần đầu hoặc áp dụng thử:}
\dotfill
\ghichu{Ghi ngày chị bắt đầu cho học sinh làm bài trên hệ thống. Nếu không nhớ
chính xác, lấy mốc đầu học kỳ (ví dụ 05/9/2025).}

\bigskip
\noindent\textbf{5. Mô tả bản chất của sáng kiến}

\subsection*{5.1. Tình trạng giải pháp đã biết}

Việc biên soạn đề kiểm tra môn Toán ở trường phổ thông hiện nay chủ yếu được
thực hiện thủ công: giáo viên sưu tầm câu hỏi từ nhiều nguồn rời rạc, gõ lại
bằng Word hoặc \LaTeX, tự cân đối ma trận rồi tự tạo các mã đề bằng cách đảo
thứ tự câu và phương án. Cách làm này bộc lộ bốn hạn chế: mất nhiều thời gian;
khó bảo đảm đúng ma trận theo bốn mức độ nhận thức; các mã đề chỉ khác nhau về
thứ tự nên thực chất vẫn là một đề; và kho câu hỏi không được mã hoá thống nhất
nên gần như không tái sử dụng được ở các năm học sau.

Về phía học sinh, nhu cầu tự luyện tập không được đáp ứng đúng cách: các em
chỉ tìm được đề trôi nổi trên mạng, thường lệch phạm vi chương trình hoặc sai
lời giải; làm xong phải chờ đến tiết sau mới biết đúng sai nên phản hồi mất
tính thời sự; và không có cách nào luyện lại nhiều lần cùng một dạng toán với
số liệu khác nhau.

Một số phần mềm tạo đề thương mại đã có trên thị trường nhưng hoặc thu phí,
hoặc chỉ hoán vị câu hỏi trong một kho có sẵn của nhà cung cấp, không cho phép
giáo viên đưa ngân hàng câu hỏi của chính mình vào, và chưa cập nhật kịp cấu
trúc định dạng đề mới gồm bốn dạng câu hỏi (trắc nghiệm nhiều lựa chọn,
đúng/sai, trả lời ngắn, tự luận) theo Quyết định số 764/QĐ-BGDĐT ngày
08/3/2024 của Bộ Giáo dục và Đào tạo. (Phân tích chi tiết ở Mục II.2.)

\subsection*{5.2. Nội dung giải pháp đề nghị công nhận là sáng kiến}

\textbf{a) Mục đích của giải pháp}

Xây dựng một hệ thống ngân hàng đề môn Toán chạy trên nền web dùng chung cho
cả giáo viên và học sinh. Về phía học sinh: chủ động tạo đề luyện tập đúng phạm
vi mình cần bằng một câu yêu cầu tiếng Việt, làm bài trực tuyến và nhận điểm
cùng đáp án ngay. Về phía giáo viên: có đề kiểm tra đúng ma trận chỉ sau vài
phút, mỗi lần tạo cho ra một đề thực sự khác nhau về số liệu, giao đề và trả
điểm về lớp học Google Classroom, đồng thời theo dõi được điểm trung bình và
những bài học sinh hay sai.

\textbf{b) Tính mới của giải pháp}

\begin{enumerate}[label=(\arabic*), leftmargin=1.2cm, itemsep=2pt]
  \item Câu hỏi không được lưu ở dạng văn bản tĩnh mà được lưu dưới dạng
        \emph{hàm sinh} viết bằng Python. Mỗi lần gọi, hàm sinh ra một câu hỏi
        mới với bộ số liệu mới và lời giải tương ứng, nên từ một câu gốc có thể
        tạo ra hàng trăm biến thể cùng dạng, cùng độ khó.
  \item Toàn bộ kho câu hỏi được mã hoá theo một chuẩn định danh thống nhất
        gắn với lớp -- chương -- bài -- mức độ -- dạng câu, nhờ đó máy có thể
        tự chọn câu theo ma trận mà không cần giáo viên chỉ định từng câu.
  \item Cùng một ngân hàng câu hỏi phục vụ hai đối tượng: dữ liệu học sinh
        làm bài quay trở lại thành thông tin thống kê cho giáo viên, tạo thành
        một vòng khép kín giữa dạy và học.
  \item Thuật toán phân bổ câu hỏi bám theo số tiết của từng bài trong phân
        phối chương trình, đồng thời xử lý riêng dạng câu đúng/sai theo quy
        ước một câu đúng/sai tương đương bốn ý ở bốn mức độ khác nhau.
  \item Trí tuệ nhân tạo chỉ được dùng ở khâu hiểu yêu cầu bằng ngôn ngữ tự
        nhiên của giáo viên; toàn bộ khâu chọn câu, tính điểm và sinh đề do mã
        nguồn đảm nhiệm. Nguyên tắc ``ưu tiên mã nguồn hơn trí tuệ nhân tạo''
        này giúp kết quả ổn định, kiểm chứng được và không phụ thuộc vào việc
        mô hình có trả lời đúng hay không.
  \item Hệ thống khép kín một vòng: từ yêu cầu của giáo viên đến đề PDF, lời
        giải PDF, bài làm trực tuyến, điểm số và sổ điểm trên Google Classroom.
\end{enumerate}

\textbf{c) Cách thức thực hiện}

Trình bày chi tiết tại Mục II.3 của báo cáo, gồm sáu bước: chuẩn hoá mã câu
hỏi; viết hàm sinh câu hỏi bằng Python; xây dựng bảng ma trận đề; lập trình
thuật toán chọn câu; sinh đề và lời giải bằng \LaTeX; tổ chức làm bài, chấm
điểm và trả kết quả về Google Classroom.

\subsection*{5.3. Khả năng áp dụng của giải pháp}

Sáng kiến đã được áp dụng tại \DONVI\ cho cả giáo viên tổ Toán và học sinh,
và có thể áp dụng cho mọi tổ Toán ở trường trung học cơ sở, trung học phổ
thông. Do phần khung hệ thống tách rời
khỏi nội dung câu hỏi, các môn học khác có ngân hàng câu hỏi dạng tính toán
(Vật lí, Hoá học, Sinh học) có thể dùng lại toàn bộ phần mềm, chỉ cần bổ sung
hàm sinh câu hỏi của môn mình.

\subsection*{5.4. Hiệu quả, lợi ích thu được}

\ghichu{Mục này em sẽ dựng bảng số liệu ở phần sau; chị điền số thật rồi chép
tóm tắt hai đến ba dòng vào đây. Không nên ghi số ước lượng trong đơn.}

Trình bày chi tiết tại Mục II.5 của báo cáo.

\subsection*{5.5. Tài liệu kèm theo}

Báo cáo sáng kiến; phụ lục hình ảnh sản phẩm; địa chỉ truy cập hệ thống;
tài khoản dùng thử dành cho hội đồng chấm.

\bigskip
\noindent\textbf{6. Những thông tin cần được bảo mật:} Không có.

\noindent\textbf{7. Các điều kiện cần thiết để áp dụng sáng kiến:} Máy tính
hoặc điện thoại có kết nối Internet; giáo viên biết sử dụng máy tính ở mức cơ
bản; nhà trường có lớp học Google Classroom (không bắt buộc).

\bigskip
\noindent Tôi xin cam đoan mọi thông tin nêu trong đơn là trung thực, đúng sự
thật và hoàn toàn chịu trách nhiệm trước pháp luật.

\vspace{1cm}
\begin{flushright}
\begin{minipage}{6cm}
\centering
\textit{\DIADIEM, ngày \dots\ tháng \dots\ năm 2026}\\[2pt]
\textbf{Người nộp đơn}\\
\textit{(Ký và ghi rõ họ tên)}\\[2.5cm]
\TACGIA
\end{minipage}
\end{flushright}

\clearpage
